% ECONOMICS_PROBLEM: {"id": "C65-33", "title": "Integrable terminal data at the critical logarithmic BSDE growth rate", "jel_code": "C65", "source_id": "OP-0562"}
% FIRST_POSTED: 2026-10-09
% ATTEMPT_STATUS: unresolved
% ENTRY_KIND: research_attempt
% !TeX program = pdflatex
% Self-contained: no external input or bibliography files required.
\documentclass[11pt,a4paper]{article}
\usepackage[T1]{fontenc}
\usepackage[utf8]{inputenc}
\usepackage{lmodern}
\usepackage[margin=24mm]{geometry}
\usepackage{amsmath,amssymb,amsthm,mathtools}
\usepackage{booktabs,longtable,array,enumitem,microtype,xurl,hyperref}
\hypersetup{colorlinks=true,linkcolor=blue,citecolor=blue,urlcolor=blue,pdftitle={Checked research attempts: nine economics and stochastic-analysis problems}}
\setlength{\parindent}{0pt}
\setlength{\parskip}{3pt}
\setlength{\emergencystretch}{3em}
\widowpenalty=10000\clubpenalty=10000\displaywidowpenalty=10000
\setlist[enumerate]{leftmargin=*,itemsep=2pt,topsep=3pt}
\newcommand{\R}{\mathbb R}
\newcommand{\N}{\mathbb N}
\newcommand{\E}{\mathbb E}
\newcommand{\Pp}{\mathbb P}
\newcommand{\Q}{\mathbb Q}
\newcommand{\F}{\mathcal F}
\newcommand{\ind}{\mathbf 1}
\newcommand{\cE}{\mathcal E}
\newcommand{\Law}{\operatorname{Law}}
\newcommand{\sgn}{\operatorname{sgn}}
\newcommand{\Var}{\operatorname{Var}}
\newcommand{\dist}{\operatorname{dist}}
\newcommand{\KL}{\operatorname{KL}}
\newcommand{\TV}{\operatorname{TV}}
\newcommand{\Lip}{\operatorname{Lip}}
\DeclareMathOperator*{\esssup}{ess\,sup}
\DeclareMathOperator*{\argmax}{arg\,max}
\newtheorem{theorem}{Theorem}
\newtheorem{lemma}[theorem]{Lemma}
\newtheorem{proposition}[theorem]{Proposition}
\newtheorem{corollary}[theorem]{Corollary}
\theoremstyle{definition}
\newtheorem{example}[theorem]{Example}
\newcommand{\report}[3]{\clearpage\setcounter{theorem}{0}\renewcommand{\thetheorem}{#1.\arabic{theorem}}\renewcommand{\theHtheorem}{#1.\arabic{theorem}}\section*{#1: #2}\addcontentsline{toc}{section}{#1: #2}\textbf{Outcome:} #3.\par\textbf{Tools:} web browsing available; Python computation available. Literature checked on 9 October 2026.}
\newcommand{\stage}[2]{\subsection*{#1) #2}}
\newcommand{\unresolved}{\textbf{LABEL: UNRESOLVED.}}
\newcommand{\disproof}{\textbf{LABEL: FULL SOLUTION. SUBLABEL: COUNTEREXAMPLE/DISPROOF.}}

\begin{document}
\section*{C65-33: Existence at the logarithmic \texorpdfstring{$L^1$}{L1} endpoint}
\textbf{Outcome:} UNRESOLVED; constructive $L^2$ existence and a necessary identity proved.\par\textbf{Tools:} web browsing available; Python computation available. Literature checked on 9 October 2026.
\subsection*{1) FORMAL RESTATEMENT}
\paragraph{Brownian conventions used in this report.}
Fix $T>0$ and an integer $d\ge1$. The probability space carries a $d$-dimensional Brownian motion $B$ with its completed, right-continuous natural filtration; $\F=\F_T$ and $\F_0$ is trivial modulo null sets. A scalar solution consists of a continuous adapted real process $Y$ and a predictable $\R^d$-valued process $Z$ such that
\[
 \int_0^T |Z_t|^2dt+\int_0^T|g(t,Y_t,Z_t)|dt<\infty\quad\Pp\text{-a.s.},
\]
and the displayed BSDE holds simultaneously for all times outside one null set. Equality of solutions means indistinguishability of $Y$ and $dt\otimes\Pp$-a.e. equality of $Z$. A process $X$ is of class (D) when $\{X_\tau:\tau\le T\text{ a stopping time}\}$ is uniformly integrable. Logs are natural, norms Euclidean, and $\sgn(0)=0$. At $T=0$ the assertions reduce to $Y_0=\xi$, with no nontrivial $Z$ component.

\paragraph{Precisely stated probabilistic tools.}
We use It\^o's formula and its Tanaka version for continuous semimartingales, Brownian martingale representation, and the following forms of change of measure. If $\theta$ is predictable and $\E\exp(\frac12\int_0^T|\theta|^2)<\infty$, Novikov's theorem says that
$D=\cE(\int\theta\,dB)$ is a uniformly integrable martingale; under $d\Q=D_Td\Pp$, $B-\int\theta dt$ is Brownian. Here $\cE(M)=\exp(M-\langle M\rangle/2)$. Also, for a continuous Brownian martingale, the BMO hypothesis
\[
 \esssup_{\tau\le T}\E\left[\left.\int_\tau^T|\theta_t|^2dt\right|\F_\tau\right]<\infty
\]
implies that $D$ is uniformly integrable and $D_T\in L^q$ for some $q>1$; this is the continuous-martingale BMO exponential theorem \cite{Kazamaki}. Each use below verifies the relevant integrability hypothesis. A local submartingale or supermartingale of class (D) is a true submartingale or supermartingale: apply the localized conditional inequalities and pass to the limit using uniform integrability. A local martingale of class (D) therefore equals the conditional expectation of its terminal value.

For every $\xi\in L^1(\F_T)$, does there exist a solution with $Y$ of class (D) to
\[
 Y_t=\xi+\int_t^T g(Z_s)ds-\int_t^TZ_s\,dB_s,
 \qquad g(z)=\frac{|z|}{\sqrt{\log(e+|z|)}}?
\]
The coefficient, logarithm, and exponent $-1/2$ are fixed. Terminal values may have either sign and need not have any $L\log^q L$ moment. The desired class-(D) condition is on $Y$, not $|Y|^p$ or $e^{|Y|}$. A construction of some local solution would not suffice. The main stress point is a very heavy but integrable positive terminal tail.

\subsection*{2) QUICK LITERATURE/CONTEXT CHECK}
This is Problem 5.1 of Fan--Hu--Tang \cite{FHT}. That source explicitly distinguishes the critical generator from its slower-growth and positive-log-moment neighboring regimes. The argument below proves an elementary square-integrable subclass without importing an endpoint existence assertion from a stronger theorem.

\subsection*{3) ATTACK PLAN}
\textbf{Proof routes.} Use the generator's Lipschitz constant to solve bounded or square-integrable approximations; then seek uniform class-(D) control for truncations. A direct integrable martingale formulation is another route.

\textbf{Disproof routes.} Test constants and affine Brownian terminals first; then construct an explicit $L^1$ terminal outside every positive log-moment class and examine whether the necessary generator-integrability identity fails. No such failure is established here.

\textbf{Landscape and tools.} This is endpoint stochastic existence. Radial differentiation identifies the exact Lipschitz constant; Brownian representation constructs the martingale part; a contraction on short time intervals proves existence; Doob's $L^2$ inequality gives class (D); monotone convergence controls a nonnegative generator integral; de la Vall\'ee Poussin-type uniform integrability is the missing compactness issue; and a probability-integral transform provides explicit terminal tails.

\subsection*{4) WORK}
\begin{lemma}[Global Lipschitz bound]
The generator is nonnegative, $g(0)=0$, $g(z)\le|z|$, and
$|g(z)-g(w)|\le|z-w|$ for every $z,w$.
\end{lemma}
\begin{proof}
For $f(r)=r/\sqrt{\log(e+r)}$, put $L=\log(e+r)\ge1$. Direct differentiation gives
\[
 f'(r)=L^{-1/2}\left(1-\frac{r}{2(e+r)L}\right),\qquad r\ge0.
\]
The term subtracted in parentheses lies in $[0,1/2)$, so $0<f'(r)\le1$. The mean value theorem and $||z|-|w||\le|z-w|$ yield the claim. The remaining assertions follow from $L\ge1$.
\end{proof}

\begin{theorem}[Constructive existence for $\xi\in L^2$]
For every square-integrable terminal value, a solution exists with
$\E\sup_{t\le T}|Y_t|^2<\infty$ and $\E\int_0^T|Z_t|^2dt<\infty$. In particular $Y$ is of class (D).
\end{theorem}
\begin{proof}
First work on $[a,b]$ with length $h=b-a<1$ and terminal value $\eta\in L^2(\F_b)$. Let $\mathcal H^2[a,b]$ be the Hilbert space of predictable integrands with norm
$\|z\|^2=\E\int_a^b|z_t|^2dt$, identifying integrands equal a.e. For $z\in\mathcal H^2$, the variable
$H=\eta+\int_a^b g(z_t)dt$ is square-integrable by Cauchy--Schwarz and $g(z)\le|z|$. Represent the martingale $M_t=\E[H\mid\F_t]$ on $[a,b]$ as
$M_t=M_a+\int_a^t\widehat Z_s\,dB_s$, and define $\mathcal T(z)=\widehat Z$.

For two inputs, It\^o isometry and conditional centering give
\begin{align*}
 \|\mathcal T(z)-\mathcal T(z')\|^2
 &=\E\left|\Delta H-\E[\Delta H\mid\F_a]\right|^2\\
 &\le\E|\Delta H|^2
 \le h\,\E\int_a^b|g(z_t)-g(z'_t)|^2dt
 \le h\|z-z'\|^2.
\end{align*}
Starting from zero and iterating $\mathcal T$ gives a Cauchy sequence because successive distances are bounded by a geometric sequence with ratio $\sqrt h<1$. Completeness gives a limit $Z$; continuity of $\mathcal T$ makes it a fixed point. Define
$Y_t=M_t-\int_a^tg(Z_s)ds$. Subtracting the equation at $b$ verifies the BSDE. Moreover Doob's inequality gives
$\E\sup_{a\le t\le b}|M_t|^2\le4\E|H|^2$, and
$\sup_{a\le t\le b}|\int_a^tg(Z_s)ds|^2\le h\int_a^b|Z_s|^2ds$.
These yield the asserted bounds on this interval.

Partition $[0,T]$ into finitely many intervals of length less than one. Solve the last interval, then use its square-integrable left endpoint as terminal value on the preceding interval, and continue backward. The finitely many estimates ensure global square integrability; the pieces agree at endpoints and concatenate to a continuous solution. Finally
$|Y_\tau|^2\le\sup_{t\le T}|Y_t|^2$ for every stopping time. This integrable dominating variable makes $|Y|^2$, and therefore $Y$, of class (D).
\end{proof}

\begin{proposition}[A necessary integrability identity at the endpoint]
Every class-(D) solution with $\xi\in L^1$ must satisfy
\[
 \E\int_0^Tg(Z_s)ds=Y_0-\E\xi<\infty,
 \qquad
 Y_t=\E\left[\left.\xi+\int_t^Tg(Z_s)ds\right|\F_t\right].
\]
In particular, if $\xi\ge0$, then $Y_t\ge0$ for every $t$ a.s.
\end{proposition}
\begin{proof}
Let $A_t=\int_0^tg(Z_s)ds$, an increasing finite process. Choose stopping times $\tau_m\uparrow T$ that localize $\int Z\,dB$ and bound $A$. Taking expectations in $Y_{\tau_m}=Y_0-A_{\tau_m}+\int_0^{\tau_m}Z\,dB$ gives
$\E A_{\tau_m}=Y_0-\E Y_{\tau_m}$. Continuity and class (D) imply $Y_{\tau_m}\to\xi$ in $L^1$. Monotone convergence gives the first identity and $A_T\in L^1$.

The process $M=Y+A$ is a local martingale. The stopping family of $A$ is dominated by $A_T$, so $M$ is of class (D) and is the martingale closed by $\xi+A_T$. Therefore $Y_t=\E[\xi+A_T\mid\F_t]-A_t$, proving the second identity. Nonnegativity follows from that formula, first at each time and then simultaneously by continuity.
\end{proof}

\begin{example}[A concrete terminal value beyond all positive log moments]
Let $U=\Phi(B_T^{(1)}/\sqrt T)$ and, on $\{U<e^{-e^2}\}$, put $r=\log(1/U)$. Define
\[
 \xi=\begin{cases}
 \displaystyle\frac1{U\,r(\log r)^2},&U<e^{-e^2},\\
 0,&U\ge e^{-e^2}.
 \end{cases}
\]
This is a finite nonnegative $\F_T$-measurable variable a.s. Its first moment is exactly
\[
 \E\xi=\int_{e^2}^\infty\frac{dr}{r(\log r)^2}=\frac12.
\]
For large $r$, $\log\xi=r-\log r-2\log\log r\ge r/2$. Consequently, for every $q>0$,
\[
 \E\bigl[\xi(\log(e+\xi))^q\bigr]
 \ge 2^{-q}\int_{R_q}^\infty\frac{r^{q-1}}{(\log r)^2}dr=\infty.
\]
To justify the divergence, for sufficiently large $r$ one has $(\log r)^2\le r^{q/2}$, and the lower integrand is then at least $r^{q/2-1}$, whose integral diverges. This is a difficult admissible terminal datum, not a nonexistence counterexample.
\end{example}

\subsection*{5) VERIFICATION}
\textbf{Small cases.} If $\xi=c$ is constant, $Y=c,Z=0$ solves the equation. If $\xi=c+z_0\cdot B_T$, the explicit solution is
$Y_t=c+z_0\cdot B_t+(T-t)g(z_0)$, $Z=z_0$; its square-integrable supremum makes it admissible. Thus neither zero data nor affine Gaussian data expose the endpoint obstruction.

\textbf{Computation.} On a logarithmic grid from $10^{-12}$ to $10^{12}$, supplemented by zero, the script finds $0.186797\le f'(r)\le1$. The integral defining the heavy terminal's first moment equals $1/2$ analytically. Pseudocode for an approximation search is: truncate $\xi$ at level $m$; solve by the proved short-interval Picard map; estimate stopping-time tails of $Y^{(m)}$ and integrated generator mass. The supplied numerical checks do not implement a stochastic solver or establish uniform integrability of these truncations.

\textbf{Adversarial review.} Global Lipschitz continuity alone did not justify passing from $L^2$ to $L^1$: the short-interval norm is an $L^2$ norm. The necessary identity used both nonnegativity of $g$ and class (D), not merely local integrability. A bounded sequence of first moments is not substituted for uniform integrability. The heavy-tail example is explicitly not called a disproof.

\subsection*{6) FINAL}
\unresolved

\textbf{Strongest checked partial result.} Every $L^2$ terminal datum has a class-(D) solution by a direct constructive proof. Every possible class-(D) endpoint solution must satisfy the finite expected generator-integral identity above. An explicit integrable terminal outside every $L\log^q L$, $q>0$, has been provided.

\textbf{Exact first gap in the truncation route.} For a fixed arbitrary nonnegative $\xi\in L^1$, prove that the entire family
$\{Y^{(m)}_\tau:m\ge1,\ \tau\le T\}$ of the square-integrable solutions with terminal $\xi\wedge m$ is uniformly integrable. This assertion has not been proved here. Even after it, convergence of $Z^{(m)}$ and the nonlinear drift, and treatment of signed data, would still require justification.

\textbf{Next three moves.} (1) Construct a data-dependent convex uniform-integrability test function adapted to the exact $r/\sqrt{\log(e+r)}$ growth. (2) Analyze generator-mass concentration for the displayed terminal using localized martingale representations. (3) Prove a stability theorem under a verified class-(D) bound and uniform control of $\int g(Z^{(m)})$, rather than assuming that convergence of $Y$ implies convergence of $Z$.

\textbf{Likely minimal counterexample, if any.} It would need genuinely unbounded, non-$L^2$ data. A positive tail such as the displayed one is a concrete candidate to test, but this report proves no obstruction for it and does not infer nonexistence merely from the lack of log moments.

\clearpage
\begin{thebibliography}{99}
\bibitem{Kazamaki}
Norihiko Kazamaki. \emph{Continuous Exponential Martingales and BMO}. Lecture Notes in Mathematics 1579, Springer, 1994. DOI: 10.1007/BFb0073585. Continuous BMO stochastic exponentials and reverse H\"older inequalities. \url{https://doi.org/10.1007/BFb0073585}.

\bibitem{FHT}
Shengjun Fan, Ying Hu and Shanjian Tang. \emph{1D nonlinear backward stochastic differential equations: a unified theory and applications}. arXiv:2309.06233v2, 11 February 2026. In particular Problems 5.1--5.5, PDF pp.\ 32--33; for C65-45 see (A1), (4.13)--(4.14), PDF p.\ 31. \url{https://arxiv.org/pdf/2309.06233v2}.

\end{thebibliography}
\end{document}
